% paper_rtsc_impact.tex
% From R_tau to the Second Electrical Revolution
% Author: Sheng-Kai Huang
% Compile: pdflatex paper_rtsc_impact.tex && pdflatex paper_rtsc_impact.tex
% Target: ~5 pages, Physical Review Letters style

\documentclass[prl,twocolumn,superscriptaddress,longbibliography,floatfix]{revtex4-2}

\usepackage{amsmath}
\usepackage{amssymb}
\usepackage{amsfonts}
\usepackage{bm}
\usepackage{hyperref}
\usepackage{booktabs}
\usepackage{graphicx}
\usepackage{xcolor}

% Shared macros (Paper 1 conventions)
\newcommand{\kB}{k_{\mathrm{B}}}
\newcommand{\Tr}{\mathrm{Tr}}
\newcommand{\Petz}{\mathcal{R}}
\newcommand{\chan}{\mathcal{N}}
\newcommand{\ket}[1]{|#1\rangle}
\newcommand{\bra}[1]{\langle#1|}
\newcommand{\braket}[2]{\langle#1|#2\rangle}
\newcommand{\abs}[1]{\left|#1\right|}
\newcommand{\tauasym}{\tau}
% Paper-specific macros
\newcommand{\taueff}{\tauasym_{\mathrm{eff}}}
\newcommand{\taudec}{\tauasym_{\mathrm{ph}}}
\newcommand{\Icoop}{I_{\mathrm{Cooper}}}
\newcommand{\Ipair}{I_{\mathrm{pair}}}
\newcommand{\Rtau}{R_{\tauasym}}
\newcommand{\Tc}{T_{\mathrm{c}}}
\newcommand{\oD}{\omega_{\mathrm{D}}}
\newcommand{\EF}{E_{\mathrm{F}}}
\newcommand{\ThetaD}{\Theta_{\mathrm{D}}}

\begin{document}

\title{From $R_\tau$ to the Second Electrical Revolution:\\
Timeline and Technological Implications of Room-Temperature Superconductivity}

\author{Sheng-Kai Huang}
\email{akai@fawstudio.com}
\affiliation{Independent Researcher}

\date{\today}

\begin{abstract}
We combine the $\Rtau$ screening criterion for
superconductivity---derived from the Petz recovery
framework---with current experimental progress to construct
a quantitative timeline for room-temperature superconductivity
(RTSC).  The analysis identifies four milestones:
$0\,^\circ$C superconductivity at ambient pressure (2027--2030),
room-temperature ($20\,^\circ$C) new materials (2030--2035),
first commercial products (2035--2045), and systemic
technological transformation (2040--2060).  For each
RTSC-enabled technology---lossless power transmission, compact
MRI, nuclear fusion, superconducting energy storage,
frictionless transport, and low-power computing---we provide
quantitative impact estimates grounded in physics.  Global
transmission losses of $\sim 2000\,$TWh/year would be
eliminated, saving $\sim \$200\,$B/year and
$\sim 1.5\,$Gt\,CO$_2$/year; MRI costs would decrease from
$\$1.5$--$3\,$M to $\$100$--$200\,$K per unit, enabling
universal diagnostic access; fusion reactor size and cost
could decrease by an order of magnitude.  We also delineate
what RTSC \emph{cannot} improve---quantum computing coherence,
battery chemistry, and optical communication---to temper
overenthusiasm.  The $\Rtau$ screener provides a
community-accessible tool for accelerating the search.
\end{abstract}

\maketitle

%=======================================================================
\section{Introduction}
\label{sec:intro}
%=======================================================================

The pursuit of room-temperature superconductivity (RTSC)---zero
electrical resistance at $T \geq 300\,$K and ambient
pressure---spans more than a century since Onnes's discovery of
superconductivity in mercury at $4.2\,$K~\cite{Onnes1911}.
Progress accelerated dramatically in recent years:
the compressed hydride LaH$_{10}$ achieves
$\Tc = 260\,$K $= -13\,^\circ$C at $170\,$GPa~\cite{Somayazulu2019,Drozdov2019},
while Deng and Chu demonstrated that pressure-quenched
HgBaCaCuO retains $\Tc = 151\,$K at ambient
pressure~\cite{DengChu2026}---the highest ambient-pressure
$\Tc$ ever recorded.

Despite this empirical momentum, the field has lacked a
\emph{quantitative} framework for predicting which materials
can achieve RTSC and on what timescale.  The $\Rtau$ screening
criterion~\cite{HuangRTSC2026}, derived from the Petz recovery
formalism~\cite{Huang2026,Petz1988}, provides exactly this:
\begin{equation}\label{eq:Rtau}
  \Rtau(T) \equiv \frac{\Ipair(T)}{\taudec(T)} \geq 1
  \quad\Longrightarrow\quad \text{superconductivity at $T$}\,,
\end{equation}
where $\Ipair$ is the total pairing information (summed over
all channels) and $\taudec$ is the phonon-induced decoherence.
The criterion is mechanism-independent and computable from
materials database quantities~\cite{HuangRTSC2026}.

In this paper, we leverage the $\Rtau$ framework together with
current experimental data to construct a timeline for RTSC
milestones (Sec.~\ref{sec:timeline}), provide quantitative
impact estimates for six key technologies
(Sec.~\ref{sec:impact}), delineate what RTSC cannot achieve
(Sec.~\ref{sec:cannot}), and outline a research roadmap
(Sec.~\ref{sec:roadmap}).

%=======================================================================
\section{Timeline predictions based on $\Rtau$}
\label{sec:timeline}
%=======================================================================

The following predictions are grounded in two inputs:
(i)~$\Rtau$ values computed for known and predicted materials
in Ref.~\cite{HuangRTSC2026}, and (ii)~the demonstrated
success of pressure-quench techniques~\cite{DengChu2026}.
Confidence levels reflect the number of independent assumptions
required.

\subsection{Prediction 1 (2027--2030): $0\,^\circ$C at ambient pressure}

LaH$_{10}$ at $170\,$GPa has $\Rtau = 3.47$, far exceeding
unity~\cite{HuangRTSC2026}.  Its $\Tc = 260\,$K $= -13\,^\circ$C
is already within $13\,$K of the $0\,^\circ$C threshold.  The
critical question is whether the $Fm\bar{3}m$ clathrate
structure can be metastably retained at ambient pressure.

The pressure-quench approach demonstrated by Deng and
Chu~\cite{DengChu2026} for HgBaCaCuO---a cuprate with very
different chemistry---establishes the viability of kinetic
trapping.  Applying the same technique to LaH$_{10}$:
rapid decompression at $T < 20\,$K exploits the negligible
hydrogen diffusivity ($D_H \sim 10^{-30}\,$cm$^2$/s) to
freeze the clathrate cage.

Even with partial structural decomposition reducing $\Tc$ by
30\%, the $\Rtau$ analysis gives
\begin{equation}\label{eq:quench}
  \Rtau^{\mathrm{quench}} \approx 0.7 \times 3.47 = 2.43 \gg 1\,,
\end{equation}
yielding $\Tc \sim 180$--$230\,$K $\approx -90$ to
$-40\,^\circ$C in the pessimistic case, and
$\Tc \sim 240$--$260\,$K $\approx -30$ to $-13\,^\circ$C
if the structure is largely preserved.  Encapsulation in
amorphous boron nitride could suppress grain-boundary
hydrogen loss~\cite{HuangRTSC2026}.

\emph{Confidence: HIGH.}  The physics (large $\Rtau$ surplus)
and the technique (pressure quench) are both demonstrated;
the remaining challenge is engineering.

\subsection{Prediction 2 (2030--2035): Room temperature at ambient pressure}

For true RTSC ($\Tc > 293\,$K at 1\,atm), the $\Rtau$
framework identifies three material families that could
achieve $\Rtau(300\,\mathrm{K}) \geq 1$ through multi-channel
pairing~\cite{HuangRTSC2026}:

\emph{(a) Clathrate hydrides} (e.g., LaBeH$_8$-type, CaH$_6$
variants).---Computationally predicted $\Tc > 300\,$K under
pressure.  Materials with Debye temperature $\ThetaD > 1200\,$K
and electron-phonon coupling $\lambda > 2.5$ achieve
$\Rtau > 1$ at $300\,$K via a single phonon channel.  The
challenge: retaining these parameters at ambient pressure.

\emph{(b) Nickelate hydrides} (e.g., La$_3$Ni$_2$O$_7$+H).---
Adding hydrogen to the bilayer nickelate introduces a fourth
pairing channel.  With three coherent channels each contributing
$\Rtau^{(i)} \sim 0.4$, the multi-channel sum
$\Rtau^{\mathrm{multi}} = \sum_i \Rtau^{(i)} \sim 1.2$
exceeds unity~\cite{HuangRTSC2026}.

\emph{(c) Diamond-scaffold hydrides} (H:B:C:N).---Hydrogen
and nitrogen co-doped boron-carbon diamond combines the
highest known $\ThetaD$ ($> 2000\,$K) with multiple
dopant-induced pairing channels.

\emph{Confidence: MEDIUM.}  Requires discovery of new
ambient-pressure materials or successful metastable
stabilization of high-pressure phases.

\subsection{Prediction 3 (2035--2045): Commercial RTSC products}

Historical precedent sets the timescale from discovery to
commercialization.  YBa$_2$Cu$_3$O$_{7-\delta}$ (YBCO) was
discovered in 1987~\cite{Wu1987}; the first commercial YBCO
wire entered the market in 1996---a 9-year gap~\cite{Larbalestier2001}.
MgB$_2$ ($\Tc = 39\,$K, discovered 2001) reached commercial
wire by 2007~\cite{Tomsic2007}.  Both required solving
engineering challenges: wire drawing, contact resistance,
critical current density ($J_c > 10^4\,$A/cm$^2$ for practical
use), and upper critical field ($H_{c2} > 10\,$T for
magnets).

For RTSC, the absence of cryogenics removes the most expensive
and fragile subsystem, but the material-specific challenges
(grain-boundary weak links, anisotropy, brittleness) will still
require $\sim 10$ years of engineering development.  First
products: RTSC power cables for short-distance urban grids and
compact MRI magnets.

\emph{Confidence: MEDIUM.}  Depends on Predictions 1--2.

\subsection{Prediction 4 (2040--2060): RTSC-enabled transformation}

Systemic transformation---global lossless power grids,
mass-produced personal maglev, portable MRI in every
ambulance---requires not just an RTSC material but mature
manufacturing at scale: kilometer-length wires, reliable
joints, and $J_c > 10^5\,$A/cm$^2$ in bulk form.  This
parallels the semiconductor industry's trajectory from the
transistor (1947) to the integrated circuit (1958) to the
microprocessor (1971)---a 24-year arc.

\emph{Confidence: LOW.}  Long-range extrapolation with many
intervening factors.

\begin{table}[t]
\caption{\label{tab:timeline}%
$\Rtau$-based RTSC milestone timeline.  $\Rtau$ values and
material data from Ref.~\cite{HuangRTSC2026}.}
\begin{ruledtabular}
\begin{tabular}{lccc}
Milestone & Period & Key $\Rtau$ & Confidence \\
\hline
$0\,^\circ$C, ambient $P$ & 2027--2030 & 2.4--3.5 & High \\
$20\,^\circ$C, ambient $P$ & 2030--2035 & $\geq 1.0$ & Medium \\
Commercial products & 2035--2045 & --- & Medium \\
Systemic transformation & 2040--2060 & --- & Low \\
\end{tabular}
\end{ruledtabular}
\end{table}

%=======================================================================
\section{Quantitative impact analysis}
\label{sec:impact}
%=======================================================================

For each RTSC-enabled technology, we provide physics-grounded
estimates.  All cost figures are in 2025 US dollars.

\subsection{Lossless power transmission}

Global electricity transmission and distribution losses total
$\sim 2000\,$TWh/year, approximately 8\% of total
generation~\cite{IEA2024}.  At a global average electricity
price of $\sim \$0.10$/kWh, this represents
$\sim \$200\,$B/year in wasted energy.  The carbon footprint
of these losses, assuming the current global generation mix
(60\% fossil), corresponds to $\sim 1.5\,$Gt\,CO$_2$/year---comparable
to the entire aviation industry.

RTSC cables would eliminate resistive losses entirely.
Residual losses in an AC system arise from dielectric
hysteresis in cable insulation and eddy currents at joints,
estimated at $< 0.1\%$~\cite{Malozemoff2012}.  The
$\Rtau$ requirement for year-round outdoor operation is
stringent:
\begin{equation}\label{eq:Tc_outdoor}
  \Tc > T_{\max}^{\mathrm{outdoor}} + \Delta T_{\mathrm{margin}}
  = 50\,^\circ\mathrm{C} + 10\,^\circ\mathrm{C}
  = 333\,\mathrm{K}\,,
\end{equation}
which requires $\Rtau(333\,\mathrm{K}) > 1$.  For indoor
cables (data centers, hospitals), $\Tc > 293\,$K suffices.
The quenched-LaH$_{10}$ scenario ($\Tc \sim 250\,$K) would be
suitable for underground cables (soil temperature $\sim 10$--$15\,^\circ$C
year-round) or refrigerated enclosures cooled to
$< -20\,^\circ$C---far cheaper than liquid nitrogen
cooling of existing HTS cables.

\emph{Near-term opportunity:} Urban power grids with
underground cable runs at constant soil temperature
($\sim 15\,^\circ$C $= 288\,$K) would benefit from any
material with $\Tc > 288\,$K.

\subsection{Compact MRI}

Current clinical MRI scanners generate $1.5$--$3\,$T fields
using NbTi wire cooled to $4.2\,$K in a liquid-helium bath.
The cryogenic subsystem (cryostat, cold head, helium
management) accounts for $\sim 40\%$ of system cost and
$> 60\%$ of maintenance expenses~\cite{Rinck2018}.  A typical
$1.5\,$T clinical MRI costs $\$1.5$--$3\,$M, with annual
operating costs of $\$100$--$200\,$K (dominated by helium
and electricity for the cold head).

An RTSC MRI would eliminate the entire cryogenic subsystem:
\begin{itemize}
\item Magnet cost: $\sim \$50\,$K (wire + winding, no cryostat)
\item Electronics: $\sim \$50$--$100\,$K (gradient coils, RF, computing)
\item Total system: $\sim \$100$--$200\,$K
\item Annual operating cost: $\sim \$10\,$K (electricity only)
\end{itemize}
The $\Rtau$ requirements are more demanding than for cables:
$\Tc > 300\,$K, upper critical field $H_{c2} > 10\,$T at
operating temperature, and $J_c > 10^4\,$A/cm$^2$.  These
are achievable for materials with $\Rtau \gg 1$, since large
$\Rtau$ correlates with large gap $\Delta$ and hence large
$H_{c2} \propto \Delta^2/\EF$.

At $\$150\,$K per unit, MRI becomes feasible for rural
clinics, ambulances, and developing countries.  The addressable
market expands from $\sim 5{,}000$ units/year (current) to
$\sim 50{,}000$ units/year, creating a $\sim \$7.5\,$B/year
new market.  The public health impact---early detection of
stroke, cancer, and traumatic brain injury in settings where
CT is the only current option---is difficult to overstate.

\subsection{Nuclear fusion}

The magnetic confinement approach to fusion requires
$\sim 10$--$20\,$T fields over $\sim 10\,$m$^3$ volumes.
ITER, the international fusion experiment, uses Nb$_3$Sn
superconducting magnets cooled to $4.5\,$K at a project cost
of $\sim \$25\,$B~\cite{ITER2020}.  The cryogenic plant alone
costs $\sim \$1\,$B and consumes $\sim 30\,$MW---a significant
parasitic load for a reactor targeting $500\,$MW fusion power.

The MIT-led SPARC/ARC design~\cite{Creely2020} already
exploits high-temperature superconductors (HTS, ReBCO,
$\Tc \sim 90\,$K) to achieve $> 20\,$T fields in compact
magnets, reducing reactor size by a factor of $\sim 4$
compared to ITER.  RTSC magnets would provide two further
advantages:
\begin{enumerate}
\item \emph{No cryogenics:} eliminating the $\$1\,$B cryogenic
  plant and its $30\,$MW parasitic load.
\item \emph{Higher fields:} RTSC materials with large
  $\Rtau$ are expected to have large gaps $\Delta$ and hence
  large $H_{c2}$, potentially enabling $> 30\,$T fields that
  further shrink the plasma volume.
\end{enumerate}
The net effect: fusion power plants $\sim 10\times$ smaller
and $\sim 10\times$ cheaper than ITER-class devices.  If RTSC
materials are available by 2035, the timeline for commercial
fusion power could accelerate from $\sim 2060$+ to
$\sim 2045$---a 15-year advance with enormous implications
for decarbonization.

\subsection{Superconducting magnetic energy storage}

A persistent current in a superconducting loop stores energy
indefinitely with zero resistive loss.  Current superconducting
magnetic energy storage (SMES) systems operate at $4\,$K and
are limited to small scales ($\sim 1$--$10\,$MJ) by cryogenic
cost~\cite{Ali2010}.

RTSC SMES would eliminate cooling costs, enabling grid-scale
deployment.  A $1\,$MWh ($= 3.6\,$GJ) RTSC solenoid with
$B = 5\,$T requires a coil volume
\begin{equation}\label{eq:smes}
  V = \frac{2\mu_0 E}{B^2}
  = \frac{2 \times 4\pi \times 10^{-7} \times 3.6 \times 10^9}
  {25}
  \approx 360\;\mathrm{m}^3\,,
\end{equation}
roughly the size of a standard shipping container
($\sim 33$--$76\,$m$^3$ for 20--40\,ft) when wound as a
tightly packed toroid at $10\,$T (reducing volume by
$4\times$).  At $10\,$T and compact geometry,
$V \sim 90\,$m$^3$---two standard containers.

Unlike lithium-ion batteries, SMES has no cycle degradation,
no fire risk, and $> 95\%$ round-trip efficiency (limited
only by the power electronics interface).  RTSC SMES could
replace lithium-ion for grid-scale storage applications where
cycle life and safety are paramount.

\subsection{Transportation: frictionless maglev}

The Shanghai Transrapid maglev, operating since 2004, uses
conventional superconducting magnets (NbTi at $4.2\,$K) and
cost $\sim \$1.3\,$B for $30\,$km of
track~\cite{Shanghai2006}---$\$43\,$M/km, compared to
$\sim \$5$--$15\,$M/km for conventional high-speed rail.
The dominant cost driver is the cryogenic infrastructure
embedded in every vehicle.

RTSC maglev eliminates vehicle-side cryogenics entirely.
The track consists of passive RTSC rails carrying persistent
currents; the vehicle carries permanent magnets (or vice
versa).  Estimated infrastructure cost reduction:
60--80\%, to $\sim \$10$--$15\,$M/km---competitive with
conventional rail.

The physics is compelling: in the zero-drag limit, energy
consumption is determined solely by aerodynamic resistance.
For a maglev train at $300\,$km/h with $C_D A \approx 8\,$m$^2$
and air density $\rho_{\mathrm{air}} = 1.2\,$kg/m$^3$:
\begin{equation}\label{eq:drag}
  P_{\mathrm{aero}} = \tfrac{1}{2}\rho_{\mathrm{air}} C_D A\, v^3
  \approx 3.3\;\mathrm{MW}\,,
\end{equation}
versus $\sim 8$--$12\,$MW for wheel-on-rail trains at the
same speed (including rolling friction).  Energy savings:
$\sim 60$--$70\%$ per passenger-km.

\subsection{Computing: superconducting interconnects}

In modern data centers, $\sim 30$--$50\%$ of power consumption
is attributable to resistive losses in wiring, power
distribution, and cooling of the associated
heat~\cite{Shehabi2016}.  Global data center electricity
consumption is $\sim 400$--$500\,$TWh/year (2024), implying
$\sim 150$--$250\,$TWh/year dissipated as wiring heat.

RTSC interconnects would eliminate resistive wiring losses,
reducing data center power consumption by $\sim 30$--$50\%$.
At $\$0.08$/kWh (industrial rate), this represents
$\$12$--$20\,$B/year in savings.

\emph{Caveat:} Superconductivity improves \emph{wiring}, not
\emph{switching}.  CMOS transistor operation does not require
superconductivity; the logic gates themselves are
non-superconducting.  Josephson junction logic~\cite{Likharev1991}
could in principle replace CMOS, but its clock speeds
($\sim 100\,$GHz, already demonstrated at $4\,$K) come with
extremely low logic density.  RTSC would make Josephson logic
more accessible by removing cryogenics, but it would not
solve the density problem.

%=======================================================================
\section{What RTSC cannot do}
\label{sec:cannot}
%=======================================================================

Balanced assessment requires identifying technologies that
will \emph{not} benefit from RTSC:

\emph{(i) Quantum computers.}---Superconducting qubits
(transmons) operate at $\sim 15\,$mK, not because they
require superconducting wiring, but because qubit
coherence demands thermal excitation energy
$\kB T \ll \hbar\omega_q \sim 25\,\mu$eV ($\omega_q/2\pi
\sim 5\,$GHz).  At room temperature ($\kB T \approx 25\,$meV),
thermal occupation $\bar{n} \sim 10^3$ destroys coherence
regardless of whether the circuit is superconducting.  RTSC
would reduce cryogenic costs for \emph{classical control
electronics} but not eliminate the millikelvin stage.

\emph{(ii) Battery charging speed.}---Charging rate is limited
by electrochemical kinetics (lithium-ion diffusion in the
cathode, SEI layer impedance), not wiring resistance.  RTSC
cables to the charger would reduce $I^2 R$ heating in the
cable but not change the fundamental charge rate of the
battery itself.

\emph{(iii) Optical communication.}---Photons have zero
electrical resistance by definition.  Fiber-optic
transmission losses ($\sim 0.2\,$dB/km at $1550\,$nm) arise
from Rayleigh scattering and absorption, neither of which
involves electrical resistance.  RTSC is irrelevant to
optical networks.

\emph{(iv) Semiconductor logic speed.}---The clock speed of
modern processors is limited by RC delay in interconnects and
power dissipation in transistors, not by wire resistance alone.
While RTSC eliminates the R in RC, the capacitance C remains,
and the switching energy per transistor (set by $\sim C V^2$)
is unchanged.

%=======================================================================
\section{The $\Rtau$ roadmap}
\label{sec:roadmap}
%=======================================================================

The path from the current state of the art to RTSC-enabled
technology can be organized into three phases, each gated
by the $\Rtau$ criterion:

\emph{Phase 1 (now--2030): Pressure-quench existing materials.}---
LaH$_{10}$ ($\Rtau = 3.47$) and H$_3$S ($\Rtau = 1.89$)
are the best candidates for kinetic trapping at ambient
pressure.  The key engineering challenge is rapid, controlled
decompression at cryogenic temperatures.  Success criterion:
$\Tc > 273\,$K at $P = 1\,$atm.

\emph{Phase 2 (2030--2040): New ambient-pressure materials.}---
The three material families identified in
Ref.~\cite{HuangRTSC2026}---clathrate hydrides, nickelate
hydrides, and diamond-scaffold hydrides---are priority targets
for first-principles screening and experimental synthesis.
The open-source $\Rtau$ screener~\cite{rtau_code} enables
community-driven search across the $> 150{,}000$ entries in
the Materials Project database~\cite{Jain2013}.  Success
criterion: ambient-pressure material with $\Rtau(300\,\mathrm{K}) > 1$.

\emph{Phase 3 (2040+): Engineering at scale.}---Wire drawing,
joining, and critical current optimization for the winning
material.  Historical precedent (YBCO, MgB$_2$) suggests
$\sim 10$ years from material discovery to first commercial
products.  Success criterion: $J_c > 10^5\,$A/cm$^2$ in
kilometer-length wires.

\subsection{The role of community search}

The $\Rtau$ criterion reduces the RTSC search to a tractable
combinatorial problem: for each candidate material, compute
$\ThetaD$, $\lambda$, and the number of pairing channels from
first-principles or database values, evaluate $\Rtau(300\,\mathrm{K})$,
and rank.  The $\Rtau$ screener code~\cite{rtau_code}
implements this pipeline and is freely available.  We
encourage the materials science community---particularly
groups with access to high-throughput DFT
workflows~\cite{Curtarolo2012}---to extend the screening
beyond the 897 hydrogen-metallic compounds examined in
Ref.~\cite{HuangRTSC2026} to the full inorganic materials
space.

%=======================================================================
\section{Conclusion}
\label{sec:conclusion}
%=======================================================================

The $\Rtau$ framework, derived from quantum information
theory, provides the first quantitative criterion for
screening RTSC candidates and enables physics-grounded
timeline predictions.  The technological implications of RTSC
are vast---\$200\,B/year in energy savings, universal MRI
access, accelerated fusion, and frictionless transport---but
they require materials with $\Tc > 300\,$K at ambient
pressure, a goal that the $\Rtau$ analysis identifies as
challenging but physically achievable.

The next experimental milestone is clear:
pressure-quench LaH$_{10}$ to ambient conditions and measure
$\Tc$.  If the quenched material retains $\Tc > 200\,$K, the
second electrical revolution will be underway.

\begin{acknowledgments}
The author thanks QuTech and the Quantum Inspire team for cloud
access to the Tuna-9 processor, and the Materials Project team
for the open database.  This work was supported in part by
the $\tauasym = 1 - F$ research program.
\end{acknowledgments}

%=======================================================================
\begin{thebibliography}{30}

\bibitem{Onnes1911}
H.~K.~Onnes,
The resistance of pure mercury at helium temperatures,
Commun.\ Phys.\ Lab.\ Univ.\ Leiden \textbf{12}, 120 (1911).

\bibitem{BCS1957}
J.~Bardeen, L.~N.~Cooper, and J.~R.~Schrieffer,
Theory of superconductivity,
Phys.\ Rev.\ \textbf{108}, 1175 (1957).

\bibitem{Bednorz1986}
J.~G.~Bednorz and K.~A.~M\"uller,
Possible high $T_c$ superconductivity in the Ba--La--Cu--O system,
Z.\ Phys.\ B \textbf{64}, 189 (1986).

\bibitem{Somayazulu2019}
M.~Somayazulu \textit{et al.},
Evidence for superconductivity above 260\,K in lanthanum superhydride at megabar pressures,
Phys.\ Rev.\ Lett.\ \textbf{122}, 027001 (2019).

\bibitem{Drozdov2019}
A.~P.~Drozdov \textit{et al.},
Superconductivity at 250\,K in lanthanum hydride under high pressures,
Nature \textbf{569}, 528 (2019).

\bibitem{Drozdov2015}
A.~P.~Drozdov, M.~I.~Eremets, I.~A.~Troyan, V.~Ksenofontov, and S.~I.~Shylin,
Conventional superconductivity at 203\,kelvin at high pressures in the sulfur hydride system,
Nature \textbf{525}, 73 (2015).

\bibitem{DengChu2026}
D.~Deng and C.~W.~Chu,
Pressure-quenched HgBaCaCuO with $T_c = 151\,$K at ambient pressure,
Preprint (2026).

\bibitem{Wu1987}
M.~K.~Wu \textit{et al.},
Superconductivity at 93\,K in a new mixed-phase Y--Ba--Cu--O compound system at ambient pressure,
Phys.\ Rev.\ Lett.\ \textbf{58}, 908 (1987).

\bibitem{HuangRTSC2026}
S.-K.~Huang,
Room-temperature superconductivity as a quantum information problem:
The $R_\tau$ screening criterion and material predictions,
Preprint (2026).

\bibitem{Huang2026}
S.-K.~Huang,
Temporal asymmetry as Petz recovery failure,
Zenodo DOI: 10.5281/zenodo.18897853 (2026);
see also \url{https://github.com/akaiHuang/petz-recovery-unification}.

\bibitem{Petz1988}
D.~Petz,
Sufficiency of channels over von Neumann algebras,
Q.\ J.\ Math.\ \textbf{39}, 97 (1988).

\bibitem{Jain2013}
A.~Jain \textit{et al.},
Commentary: The Materials Project: A materials genome approach to accelerating materials innovation,
APL Mater.\ \textbf{1}, 011002 (2013).

\bibitem{Larbalestier2001}
D.~Larbalestier, A.~Gurevich, D.~M.~Feldmann, and A.~Polyanskii,
High-$T_c$ superconducting materials for electric power applications,
Nature \textbf{414}, 368 (2001).

\bibitem{Tomsic2007}
M.~Tomsic \textit{et al.},
Overview of MgB$_2$ superconductor applications,
Int.\ J.\ Appl.\ Ceram.\ Technol.\ \textbf{4}, 250 (2007).

\bibitem{IEA2024}
International Energy Agency,
\textit{World Energy Outlook 2024} (IEA, Paris, 2024).

\bibitem{Malozemoff2012}
A.~P.~Malozemoff \textit{et al.},
Progress in high temperature superconductor coated conductors
and their applications,
Supercond.\ Sci.\ Technol.\ \textbf{25}, 123001 (2012).

\bibitem{Rinck2018}
P.~A.~Rinck,
\textit{Magnetic Resonance in Medicine}, 12th ed.\
(ABW Wissenschaftsverlag, Berlin, 2018).

\bibitem{ITER2020}
ITER Organization,
\textit{ITER Research Plan within the Staged Approach}
(ITER Organization, St.\ Paul-lez-Durance, 2018).

\bibitem{Creely2020}
A.~J.~Creely \textit{et al.},
Overview of the SPARC tokamak,
J.\ Plasma Phys.\ \textbf{86}, 865860502 (2020).

\bibitem{Ali2010}
M.~H.~Ali, B.~Wu, and R.~A.~Dougal,
An overview of SMES applications in power and energy systems,
IEEE Trans.\ Sustain.\ Energy \textbf{1}, 38 (2010).

\bibitem{Shanghai2006}
R.~Thornton, T.~Clark, and E.~Perreault,
Linear synchronous motor propulsion of urban transit vehicles,
in \textit{Proceedings of the Joint Rail Conference} (ASME, 2006).

\bibitem{Shehabi2016}
A.~Shehabi \textit{et al.},
\textit{United States Data Center Energy Usage Report}
(Lawrence Berkeley National Laboratory, 2016).

\bibitem{Likharev1991}
K.~K.~Likharev and V.~K.~Semenov,
RSFQ logic/memory family: a new Josephson-junction technology
for sub-terahertz-clock-frequency digital systems,
IEEE Trans.\ Appl.\ Supercond.\ \textbf{1}, 3 (1991).

\bibitem{Gao2025}
M.~Gao \textit{et al.},
Upper bound on the critical temperature of superconductors from
first-principles phonon calculations,
Phys.\ Rev.\ B \textbf{111}, 024506 (2025).

\bibitem{Curtarolo2012}
S.~Curtarolo \textit{et al.},
AFLOW: An automatic framework for high-throughput materials discovery,
Comput.\ Mater.\ Sci.\ \textbf{58}, 218 (2012).

\bibitem{rtau_code}
S.-K.~Huang,
$R_\tau$ screener: open-source code for superconductor screening,
\url{https://github.com/akaiHuang/petz-recovery-unification} (2026).

\bibitem{Huang2026exp}
S.-K.~Huang,
Measuring temporal asymmetry on a superconducting quantum processor:
From quantum erasure to the birth of time's arrow,
Preprint (2026).

\bibitem{Schilling1993}
A.~Schilling, M.~Cantoni, J.~D.~Guo, and H.~R.~Ott,
Superconductivity above 130\,K in the Hg--Ba--Ca--Cu--O system,
Nature \textbf{363}, 56 (1993).

\end{thebibliography}

\end{document}
